% Preamble
\documentclass[11pt, aspectratio=169]{beamer}
	% Packages
	\usepackage[english]{babel}
	\usepackage{lipsum}
	\usepackage{mathptmx}
	\usepackage{xcolor}
	% Themes
	\usefonttheme[onlymath]{serif}
	\usetheme[nofirafonts]{focus}
	\renewcommand{\thempfootnote}{\arabic{mpfootnote}}
	% Cover
	\title{Microeconometrics Using Stata}
	\subtitle{LINEAR REGRESSION WITH CORRELATED ERRORS}
	\author{
		Jhon R. Ordoñez
		\footnote{\label{1} National University of San Cristóbal de Huamanga}
		\footnote{\label{2} Faculty of Economic, Administrative and Accounting Sciences}
		\footnote{\label{3} Professional School of Economics}
	}
	\date{\today}
% Body
\begin{document}
	% Cover
	\begin{frame}
		\maketitle
	\end{frame}
	% Outline
	\begin{frame}[t]
		\frametitle{Outline}
		\tableofcontents
	\end{frame}
	% Section 1
	\section{Exercise 1}
		% Slide 1
		\begin{frame}
			\frametitle{Exercise 1}
				\begin{block}{Exercise 1}
					Generate data by using the same \textbf{DGP} as that in \textcolor{red}{\textbf{section 6.3}}, and
					implement the first step of \textbf{FGLS} estimation to get the predicted
					variance \textcolor{blue}{\textbf{varu}}. Now, compare several different methods to implement
					the second step of weighted estimation. First, use \textcolor{blue}{\textbf{regress}} with the
					modifier \textcolor{blue}{\textbf{[aweight=1/varu]}}, as in the text. Second, manually
					implement this regression by generating the transformed variable
					\textcolor{blue}{\textbf{try=y/sqrt(varu)}} and regressing \textcolor{blue}{\textbf{try}} on the similarly constructed
					variables \textcolor{blue}{\textbf{trx2}}, \textcolor{blue}{\textbf{trx3}}, and \textcolor{blue}{\textbf{trone}}, using \textcolor{blue}{\textbf{regress}} with the \textcolor{blue}{\textbf{noconstant}}          
					option. Third, use \textcolor{blue}{\textbf{regress}} with \textcolor{blue}{\textbf{[pweight=1/varu]}}, and show that the
					default standard errors using \textcolor{blue}{\textbf{pweight}}s differ from those using
					\textcolor{blue}{\textbf{aweight}}s because the \textcolor{blue}{\textbf{pweight}}s default is to compute robust standard
					errors.
				\end{block}
		\end{frame}
	% Section 2
	\section{Exercise 2}
		% Slide 1
		\begin{frame}[t]
			\frametitle{Exercise 2}
			\begin{block}{Exercise 2}
				Consider the same \textbf{DGP} as that in \textcolor{red!50}{\textbf{section 6.3}}. Given this specification
				of the model, the rescaled equation
					\begin{equation}
						\dfrac{y}{w} = \beta_{1}\left(\dfrac{1}{w}\right) + \beta_{2}\left(\dfrac{x_{2}}{w}\right) + \beta_{3}\left(\dfrac{x_{3}}{w}\right) + e
					\end{equation}
				where
					\begin{equation}
						w = \sqrt{exp\left(-1 + 0.2x_{2}\right)}
					\end{equation}
				will have the error $e$, which is normally distributed and homoskedastic.
				Treat $w$ as known, and estimate this rescaled regression in Stata by 
				using \textcolor{blue}{\textbf{regress}} with the \textcolor{blue}{\textbf{noconstant}} option. 
				Compare the results with those given in \textcolor{red!50}{\textbf{section 6.3}}, where the
				weight $w$ was estimated. Is there a big difference here between the \textbf{GLS} and \textbf{FGLS} estimators?
			\end{block}
		\end{frame}
	% Section 3
	\section{Exercise 3}
		% Slide 1
		\begin{frame}
			\frametitle{Exercise 3}
			\begin{block}{Exercise 3}
				Consider the same DGP as that in \textcolor{red!50}{\textbf{section 6.3}}. Suppose that we
				incorrectly assume that 
					\begin{equation}
						u \sim N\left(0, \sigma^{2}x_{2}^{2}\right)
					\end{equation}
				Then \textbf{FGLS} estimates can be obtained by using \textcolor{blue}{\textbf{regress}} with \textcolor{blue}{\textbf{[pweight=1/x2sq]}}, where \textcolor{blue}{\textbf{x2sq=x22}} .
				How different are the estimates of $(\beta_{1}, \beta_{2}, \beta_{3})$ from the \textbf{OLS} results?
				Can you explain what has happened in terms of the consequences of
				using the wrong skedasticity function? Do the standard errors change
				much if robust standard errors are computed? Use the \textcolor{blue}{\textbf{estat hettest}}
				command to check whether the regression errors in the transformed
				model are homoskedastic.
			\end{block}
		\end{frame}
	% Section 4
	\section{Exercise 4}
		% Slide 1
		\begin{frame}
			\frametitle{Exercise 4}
			\begin{block}{Exercise 4}
				For the model and data of \textcolor{red!50}{\textbf{section 6.7}}, verify that \textcolor{blue}{\textbf{mixed}} with the \textcolor{blue}{\textbf{mle}}
				option gives the same results as \textcolor{blue}{\textbf{xtreg}}, \textcolor{blue}{\textbf{mle}}. Also, compare the results
				with those from using \textcolor{blue}{\textbf{mixed}} with the \textcolor{blue}{\textbf{reml}} option. Fit the two-way RE
				model assuming random individual and time effects, and compare
				results with those from when the time effects are allowed to be fixed
				(in which case time dummies are included as regressors).
			\end{block}
		\end{frame}
	% Section 5
	\section{Exercise 5}
		% Slide 1
		\begin{frame}
			\frametitle{Exercise 5}
			\begin{block}{Exercise 5}
				Consider the same dataset as in \textcolor{red!50}{\textbf{section 6.8}}. Repeat the analysis of
				\textcolor{red!50}{\textbf{section 6.8}} using the dependent variables \textcolor{blue}{\textbf{drugexp}} and \textcolor{blue}{\textbf{totothr}}, which
				are in levels rather than logs (so heteroskedasticity is more likely to be
				a problem). First, estimate the two equations using \textbf{OLS} with default
				standard errors and robust standard errors, and compare the standard
				errors. Second, estimate the two equations jointly using \textcolor{blue}{\textbf{sureg}}, and
				compare the estimates with those from \textbf{OLS}. Third, use the \textcolor{blue}{\textbf{bootstrap}}
				prefix to obtain robust standard errors from \textcolor{blue}{\textbf{sureg}}, and compare the
				efficiency of joint estimation with that of OLS. Hint: It is much easier to
				compare estimates across methods if the \textcolor{blue}{\textbf{estimates}} command is used;
				see \textcolor{red!50}{\textbf{section 3.5.6}}.
			\end{block}
		\end{frame}
	% Section 6
	\section{Exercise 6}
		% Slide 1
		\begin{frame}
			\frametitle{Exercise 6}
			\begin{block}{Exercise 6}
				Consider the same dataset as in \textcolor{red!50}{\textbf{section 6.9}}. Repeat the analysis of
				\textcolor{red!50}{\textbf{section 6.9}} using all observations rather than restricting the sample to
				ages between 21 and 65 years. First, declare the survey design.
				Second, compare the unweighted mean of \textcolor{blue}{\textbf{hgb}} and its standard error,
				ignoring survey design, with the weighted mean and standard error,
				allowing for all features of the survey design. Third, do a similar
				comparison for least-squares regression of \textcolor{blue}{\textbf{hgb}} on \textcolor{blue}{\textbf{age}} and \textcolor{blue}{\textbf{female}}.
				Fourth, estimate this same regression using \textcolor{blue}{\textbf{regress}} with \textcolor{blue}{\textbf{pweight}s}
				and cluster-robust standard errors, and compare with the survey
				results.
			\end{block}
		\end{frame}
	% Section 7
	\section{Exercise 7}
		% Slide 1
		\begin{frame}
			\frametitle{Exercise 7}
			\begin{block}{Exercise 7}
				Reconsider the dataset from \textcolor{red!50}{\textbf{section 6.8.3}}. Estimate the parameters of
				each equation by \textbf{OLS}. Compare these \textbf{OLS} results with the \textbf{SUR} results
				reported in \textcolor{red!50}{\textbf{section 6.8.3}}.
			\end{block}
		\end{frame}
	% References
	\begin{frame}[t,allowframebreaks]
		\frametitle{References}
		\bibliography{biblio.bib}
		% Style
		\bibliographystyle{apalike}
		% Authors
		\nocite{cameron2022-I}
		\nocite{cameron2022-II}
	\end{frame}
	% Cover
	\begin{frame}
		\maketitle
	\end{frame}
\end{document}